\chapter[Geometric Predicates and Robust Computation]{Geometric Predicates\\ and Robust Computation}
\label{ch:geometric_predicates}

Geometric predicates are exact-sign tests that determine the combinatorial
relationship between geometric objects. They are the foundation of the
point, line, triangle, quadrilateral, tetrahedral, and hexahedral algorithms
that follow. A predicate should return a discrete result such as positive,
zero, or negative, rather than a constructed geometric object.

\section{Orientation Predicate}
\label{sec:orientation_predicate}

For points $A,B,C \in \mathbb{R}^2$, define
\[
  \operatorname{orient2d}(A,B,C) =
  (B_x-A_x)(C_y-A_y) - (B_y-A_y)(C_x-A_x).
\]
The sign is positive when $C$ lies to the left of the directed line
$A\mathbin{\to}B$, negative when it lies to the right, and zero when the
three points are collinear. This test is used by convex hull, intersection,
point-location, and triangulation algorithms.

\begin{algorithm}[H]
\caption{Two-dimensional orientation test}
\label{alg:orient2d}
\begin{algorithmic}[1]
\Function{Orient2D}{$A,B,C$}
  \State $u \gets B_x-A_x$
  \State $v \gets B_y-A_y$
  \State $w \gets C_x-A_x$
  \State $z \gets C_y-A_y$
  \State \Return $u z-v w$
\EndFunction
\end{algorithmic}
\end{algorithm}

\section{Three-dimensional Orientation}
\label{sec:orientation3d_predicate}

For points $A,B,C,D \in \mathbb{R}^3$, the signed volume predicate is
\[
  \operatorname{orient3d}(A,B,C,D) =
  \det\begin{pmatrix}
    B-A\\ C-A\\ D-A
  \end{pmatrix}.
\]
Its sign identifies which side of the oriented plane $ABC$ contains $D$.
Its zero case detects coplanarity. Tetrahedralization, volume-mesh
validation, and tetrahedron--hexahedron conversion use this predicate.

\section{In-circle Predicate}
\label{sec:incircle_predicate}

Given a counter-clockwise triangle $ABC$ and a point $D$, the in-circle
predicate determines whether $D$ lies inside, on, or outside the circumcircle
of $ABC$. It is evaluated from the sign of
\[
  \operatorname{incircle}(A,B,C,D) =
  \det\begin{pmatrix}
    A_x & A_y & A_x^2+A_y^2 & 1\\
    B_x & B_y & B_x^2+B_y^2 & 1\\
    C_x & C_y & C_x^2+C_y^2 & 1\\
    D_x & D_y & D_x^2+D_y^2 & 1
  \end{pmatrix}.
\]
The test is central to Delaunay triangulation and edge-flip algorithms.

\section{Robust Evaluation}
\label{sec:robust_predicates}

Floating-point roundoff can change the sign of a nearly zero determinant and
therefore change an algorithm's topology. Implementations should evaluate
predicates with adaptive precision or exact arithmetic when the rounded
result is uncertain. Epsilon thresholds may be useful for application-level
 tolerances, but they do not by themselves provide a consistent combinatorial
result. Symbolic perturbation is another way to define deterministic answers
for coincident, collinear, coplanar, or co-circular inputs.

The predicates in this chapter have constant arithmetic cost for fixed
dimension. Their practical cost is dominated by the precision required to
certify the sign.

\section{Floating-Point Filters}
\label{sec:floating_point_filters}

\subsection{Overview}
A \index{floating-point filter}\textbf{floating-point filter} evaluates a
predicate in floating point first and certifies the sign using an \emph{a
priori} error bound; only when the rounded result is too close to zero does
it fall back to a slower, exact evaluation. Because ill-conditioned inputs
are rare in practice, filters give the robustness of exact arithmetic at
nearly floating-point speed~\cite{shewchuk1997}.

\subsection{Static and Dynamic Filters}
A \textbf{static filter} derives a fixed error bound from the magnitudes of
the input coordinates before the computation begins; it is cheap but
pessimistic. A \textbf{dynamic filter} computes the bound during the
evaluation, tightening it as intermediate results are refined, so the exact
path is entered only when it is genuinely needed~\cite{bronnimann2001}.

\subsection{Adaptive Predicates}
Shewchuk's adaptive predicates evaluate the orientation and in-circle
determinants as a sequence of increasingly accurate approximations built
from \index{expansion arithmetic}\textbf{expansion arithmetic}, stopping as
soon as the sign is certified~\cite{shewchuk1997}. The arithmetic rests on
error-free transformations such as the \emph{two-sum} and \emph{two-product}
primitives of Knuth and Priest: for IEEE~754 inputs these recover the
roundoff error exactly, so the exact sign is obtained with no explicit
multiple-precision representation~\cite{priest1991,goldberg1991}. In the
common nondegenerate case the predicates run at roughly the speed of plain
floating point.

\subsection{Complexity}
Each error-free transformation costs a small constant number of floating
operations. The adaptive stage recomputes the determinant with increasing
precision only along the path where roundoff threatens the sign; the exact
stage has cost proportional to the number of bits that must be carried.

\section{Exact Arithmetic}
\label{sec:exact_arithmetic}

\subsection{Overview}
The \index{exact geometric computation}\textbf{exact geometric computation}
(EGC) paradigm treats robustness as a non-issue by computing predicates
exactly, and only when a new geometric object is constructed does it degrade
to a certified numerical value~\cite{yap1997}. Exactness is achieved with
arbitrary-precision integers, rationals, or floating-point expansions.

\subsection{Integer and Rational Kernels}
When input coordinates are given as integers (or rationals), every predicate
sign is an exact integer computation. Fortune and Van Wyk generate
specialized, unrolled evaluation code from an expression compiler so that the
cost tracks the actual bit-length of the operands rather than paying a
general-purpose library's overhead~\cite{fortune1993,fortune1996}. The bit
length grows only by a constant factor for each predicate: a 2D
orientation of $\tau$-bit integers needs $O(\tau)$ bits, and an in-circle
test needs $O(\tau)$ as well.

\subsection{Expansions}
An \index{expansion}\textbf{expansion} represents an exact value as the
exact sum of non-overlapping floating-point components. The two-sum and
two-product identities split a single floating-point addition or product into
the rounded result and its exact error, and these identities compose into
exact sums and products of arbitrary length~\cite{priest1991,shewchuk1997}.
This is the representation used by robust implementations of the orientation
and in-circle predicates.

\subsection{Exact Geometric Computation}
Yap's EGC framework adds \emph{precision-driven} evaluation: expressions are
evaluated to successively higher precision until a requested bound is met.
It underwrites libraries such as CORE that combine rational and algebraic
number types, and it extends exact computation beyond predicates to
constructed objects like intersection points~\cite{yap1997,halperin2004}.

\section{Rational Arithmetic in Computational Geometry}
\label{sec:rational_arithmetic}

\subsection{Overview}
\index{rational arithmetic}Rational arithmetic represents every number as a
ratio of arbitrary-precision integers, so all predicates and all constructed
points remain exact. It is the natural arithmetic for algorithms whose input
is integer or rational, and it avoids the roundoff of floating point at the
price of normalizing fractions and managing large numerators.

\subsection{Application to Delaunay Triangulation}
Karasick, Lieber, and Nackman showed that a Delaunay triangulation driven by
rational arithmetic eliminates the failures caused by floating-point in-circle
tests, at a cost that is acceptable for engineering input~\cite{karasick1991}.
The in-circle predicate for rational inputs reduces to a fixed combination of
integer operations whose size is bounded by the input bit-length.

\subsection{Trade-offs}
Rational arithmetic is exact and uniform, but numerators can grow large when
intersection points are repeatedly constructed. Hybrid designs combine a
floating-point filter with a rational or integer fallback so that typical
queries never touch the expensive exact kernel~\cite{fortune1993,karasick1991}.

\section{Simulation of Simplicity}
\label{sec:simulation_of_simplicity}

\subsection{Overview}
\index{Simulation of Simplicity}\textbf{Simulation of Simplicity} (SoS) is a
general technique for coping with degenerate inputs: it simulates an
infinitesimal perturbation of the data that eliminates all coincident,
collinear, coplanar, and co-circular configurations, yet the perturbation is
never actually computed~\cite{edelsbrunner1990}. The programmer writes the
algorithm for general position only; the tie-breaking is hidden inside the
predicate implementations.

\subsection{Construction}
Each coordinate $x_i$ is conceptually replaced by $x_i(\epsilon)$, a
polynomial in an indeterminate $\epsilon$, so that the perturbed predicate
$P(x_0(\epsilon), \ldots, x_{n-1}(\epsilon))$ is a nonzero polynomial whose
sign at $\epsilon \to 0^{+}$ is the lexicographically first nonzero term.
Evaluating only that leading term yields a deterministic, consistent answer
for every degenerate input. Yap proved consistency for the corresponding
symbolic perturbation scheme and extended the technique to arbitrary
predicates~\cite{yap1988,yap1990}.

\subsection{Cost and Limitations}
For determinant predicates the perturbed test costs only a small constant
factor over the unperturbed one. SoS cannot make an inconsistent
\emph{algorithm} robust, and it ties degeneracies to a fixed symbolic order
rather than to any geometric meaning; exact arithmetic with explicit
degeneracy handling remains the alternative when the tie-break matters~\cite{edelsbrunner1990,halperin2004}.

\section{Robustness Failure Examples}
\label{sec:robustness_failures}

Floating-point predicates do not merely produce slightly wrong answers; they
can make an algorithm crash, loop forever, or output an inconsistent
structure. Kettner, Mehlhorn, Pion, Schirra, and Yap collected a suite of
small, reproducible examples --- such as a convex-hull routine whose
topology changes under rotation, and a segment intersection where roundoff
contradicts the planar subdivision --- that expose these failures and are
now used as regression tests for libraries like CGAL~\cite{kettner2008}.
These examples motivate the filter and exact-arithmetic machinery above:
certifying the sign of a predicate is the only way to guarantee a consistent
combinatorial result.

